% --- Archon physics formalization source begin ---
% archon:physics
% archon:covers PhyXMiniProblems/problem_phyx_mini_0321.lean
% archon:source-report reports/phyx_mini/problem_phyx_mini_0321.source.json

\chapter{Physics problem phyx\_mini\_0321}
\label{ch:PhyXMiniProblems_problem_phyx_mini_0321}

\paragraph{Problem source.}
\#\# Physical scenario

A sperm whale vocalizes by producing a series of clicks. Actually, the whale makes only a single sound near the front of its head to start the series. Part of that sound then emerges from the head into the water to become the first click of the series. The rest of the sound travels backward through the spermaceti sac (a body of fat), reflects from the frontal sac (an air layer), and then travels forward through the spermaceti sac. When it reaches the distal sac (another air layer) at the front of the head, some of the sound escapes into the water to form the second click, and the rest is sent back through the spermaceti sac (and ends up forming later clicks). Figure also shows a strip-chart recording of a series of clicks. A unit time interval of \textbackslash{}(1.0\textbackslash{},\textbackslash{}mathrm\{ms\}\textbackslash{}) is indicated on the chart.

\#\# Question

Assuming that the speed of sound in the spermaceti sac is \textbackslash{}(1372\textbackslash{},\textbackslash{}mathrm\{m/s\}\textbackslash{}), find the length of the spermaceti sac.

\#\# Answer choices

- A: A: 3.4 m

- B: B: 2.9 m

- C: C: 2.5 m

- D: D: 2.1 m

\#\# Figure caption (auxiliary; use the image as primary evidence)

The image is a diagram of a sperm whale's head, illustrating its internal structures involved in sound production. Key components labeled are:

1. **Distal Sac**: Located at the front of the head, marked on the left.
2. **Spermaceti Sac**: A large, central structure shown in a beige color.
3. **Frontal Sac**: Positioned towards the right, at the top of the head.

Below the whale diagram, there is a waveform graph showing the pattern of sounds over a time scale. The graph includes a label indicating a time measurement of "1.0 ms" (milliseconds) between peaks, illustrating the timing of sound pulses generated by the whale.

\#\# Dataset metadata

Sound; Physical Model Grounding Reasoning, Numerical Reasoning

\paragraph{Recorded answer/context.}
D: D: 2.1 m

\paragraph{Figure/image path.}
/root/proposal\_for\_physic/science-mango/phyx\_mini\_run/phyx\_data/test\_image/321.png

\paragraph{Formalization target.}
create a compiling Lean file with sorry bodies at `PhyXMiniProblems/problem\_phyx\_mini\_0321.lean`. The Lean declarations must preserve the physical quantities, dimensions or dimensional roles, figure labels, governing-law hypotheses, and final relation expressed by this problem.
Use Mathlib/Physlib names found through LeanExplore where available. If a physics API is missing, introduce faithful local abstractions rather than scalar placeholder aliases.

\begin{theorem}[Physics formalization target]
\label{thm:physics:phyx_mini_0321:target}
\lean{PhyXMiniProblems.ProblemPhyXMini0321.spermacetiSacLength_matches_recordedAnswerD}
\uses{def:physics:phyx-mini-0321:phyxminiproblems-problemphyxmini0321-lengthinmeters, def:physics:phyx-mini-0321:phyxminiproblems-problemphyxmini0321-spermwhaleclicksetup, def:physics:phyx-mini-0321:phyxminiproblems-problemphyxmini0321-matchesproblemandfigure, def:physics:phyx-mini-0321:phyxminiproblems-problemphyxmini0321-hasphysicalparameters, def:physics:phyx-mini-0321:phyxminiproblems-problemphyxmini0321-satisfiesechoroundtriplaws, def:physics:phyx-mini-0321:phyxminiproblems-problemphyxmini0321-matchesanswerchoice, def:physics:phyx-mini-0321:phyxminiproblems-problemphyxmini0321-recordedanswerchoice}
The assigned autoformalize agent should translate this physics problem into Lean declarations in the covered file, with theorem and lemma proof bodies written as `by sorry`.
\end{theorem}
\begin{proof}
This is an autoformalization task, not a proof task. Produce statements that can later be proved by the physics prover without weakening the physical model.
\end{proof}

% --- Archon declaration topology begin ---
\section{Declaration topology}
\begin{definition}[Acoustic Length]
  \label{def:physics:phyx-mini-0321:phyxminiproblems-problemphyxmini0321-acousticlength}
  \lean{PhyXMiniProblems.ProblemPhyXMini0321.AcousticLength}
  A nonnegative physical length, independent of the unit used to read it
\end{definition}
\begin{proof}
  This is the declared model definition of Acoustic Length.
\end{proof}
\begin{definition}[Acoustic Duration]
  \label{def:physics:phyx-mini-0321:phyxminiproblems-problemphyxmini0321-acousticduration}
  \lean{PhyXMiniProblems.ProblemPhyXMini0321.AcousticDuration}
  A nonnegative physical elapsed time, independent of the readout unit
\end{definition}
\begin{proof}
  This is the declared model definition of Acoustic Duration.
\end{proof}
\begin{definition}[length Readout]
  \label{def:physics:phyx-mini-0321:phyxminiproblems-problemphyxmini0321-lengthreadout}
  \lean{PhyXMiniProblems.ProblemPhyXMini0321.lengthReadout}
  \uses{def:physics:phyx-mini-0321:phyxminiproblems-problemphyxmini0321-acousticlength}
  Read a physical length in a selected length unit
\end{definition}
\begin{proof}
  This is the declared model definition of length Readout.
\end{proof}
\begin{definition}[duration Readout]
  \label{def:physics:phyx-mini-0321:phyxminiproblems-problemphyxmini0321-durationreadout}
  \lean{PhyXMiniProblems.ProblemPhyXMini0321.durationReadout}
  \uses{def:physics:phyx-mini-0321:phyxminiproblems-problemphyxmini0321-acousticduration}
  Read a physical duration in a selected time unit
\end{definition}
\begin{proof}
  This is the declared model definition of duration Readout.
\end{proof}
\begin{definition}[speed Readout]
  \label{def:physics:phyx-mini-0321:phyxminiproblems-problemphyxmini0321-speedreadout}
  \lean{PhyXMiniProblems.ProblemPhyXMini0321.speedReadout}
  Read a physical speed in compatible selected length and time units
\end{definition}
\begin{proof}
  This is the declared model definition of speed Readout.
\end{proof}
\begin{definition}[length In Meters]
  \label{def:physics:phyx-mini-0321:phyxminiproblems-problemphyxmini0321-lengthinmeters}
  \lean{PhyXMiniProblems.ProblemPhyXMini0321.lengthInMeters}
  \uses{def:physics:phyx-mini-0321:phyxminiproblems-problemphyxmini0321-acousticlength, def:physics:phyx-mini-0321:phyxminiproblems-problemphyxmini0321-lengthreadout}
  Meter readout used for sac and path lengths in the problem
\end{definition}
\begin{proof}
  This is the declared model definition of length In Meters.
\end{proof}
\begin{definition}[duration In Seconds]
  \label{def:physics:phyx-mini-0321:phyxminiproblems-problemphyxmini0321-durationinseconds}
  \lean{PhyXMiniProblems.ProblemPhyXMini0321.durationInSeconds}
  \uses{def:physics:phyx-mini-0321:phyxminiproblems-problemphyxmini0321-acousticduration, def:physics:phyx-mini-0321:phyxminiproblems-problemphyxmini0321-durationreadout}
  Second readout used for the click-chart timing data
\end{definition}
\begin{proof}
  This is the declared model definition of duration In Seconds.
\end{proof}
\begin{definition}[speed In Meters Per Second]
  \label{def:physics:phyx-mini-0321:phyxminiproblems-problemphyxmini0321-speedinmeterspersecond}
  \lean{PhyXMiniProblems.ProblemPhyXMini0321.speedInMetersPerSecond}
  \uses{def:physics:phyx-mini-0321:phyxminiproblems-problemphyxmini0321-speedreadout}
  Meters-per-second readout used for the stated sound speed
\end{definition}
\begin{proof}
  This is the declared model definition of speed In Meters Per Second.
\end{proof}
\begin{definition}[Sac Boundary]
  \label{def:physics:phyx-mini-0321:phyxminiproblems-problemphyxmini0321-sacboundary}
  \lean{PhyXMiniProblems.ProblemPhyXMini0321.SacBoundary}
  The two air-layer boundaries labeled in the whale-head diagram
\end{definition}
\begin{proof}
  This is the declared model definition of Sac Boundary.
\end{proof}
\begin{definition}[Propagation Medium]
  \label{def:physics:phyx-mini-0321:phyxminiproblems-problemphyxmini0321-propagationmedium}
  \lean{PhyXMiniProblems.ProblemPhyXMini0321.PropagationMedium}
  The propagation medium labeled Spermaceti sac in the diagram
\end{definition}
\begin{proof}
  This is the declared model definition of Propagation Medium.
\end{proof}
\begin{definition}[Echo Route]
  \label{def:physics:phyx-mini-0321:phyxminiproblems-problemphyxmini0321-echoroute}
  \lean{PhyXMiniProblems.ProblemPhyXMini0321.EchoRoute}
  \uses{def:physics:phyx-mini-0321:phyxminiproblems-problemphyxmini0321-sacboundary, def:physics:phyx-mini-0321:phyxminiproblems-problemphyxmini0321-propagationmedium}
  A sound route through a medium from a boundary to a reflector and back
\end{definition}
\begin{proof}
  This is the declared model definition of Echo Route.
\end{proof}
\begin{definition}[Click Mechanism]
  \label{def:physics:phyx-mini-0321:phyxminiproblems-problemphyxmini0321-clickmechanism}
  \lean{PhyXMiniProblems.ProblemPhyXMini0321.ClickMechanism}
  \uses{def:physics:phyx-mini-0321:phyxminiproblems-problemphyxmini0321-sacboundary, def:physics:phyx-mini-0321:phyxminiproblems-problemphyxmini0321-echoroute}
  How a particular chart click is formed
\end{definition}
\begin{proof}
  This is the declared model definition of Click Mechanism.
\end{proof}
\begin{definition}[Sperm Whale Click Setup]
  \label{def:physics:phyx-mini-0321:phyxminiproblems-problemphyxmini0321-spermwhaleclicksetup}
  \lean{PhyXMiniProblems.ProblemPhyXMini0321.SpermWhaleClickSetup}
  \uses{def:physics:phyx-mini-0321:phyxminiproblems-problemphyxmini0321-acousticlength, def:physics:phyx-mini-0321:phyxminiproblems-problemphyxmini0321-acousticduration, def:physics:phyx-mini-0321:phyxminiproblems-problemphyxmini0321-echoroute, def:physics:phyx-mini-0321:phyxminiproblems-problemphyxmini0321-clickmechanism}
  The physical quantities and chart roles for the sperm-whale experiment. spermacetiSacLength is the requested unknown and has no numerical value in this structure. echoRoundTripDistance is kept as a separate physical length so geometry and constant-speed travel can be stated as distinct laws
\end{definition}
\begin{proof}
  This is the declared model definition of Sperm Whale Click Setup.
\end{proof}
\begin{definition}[Matches Problem And Figure]
  \label{def:physics:phyx-mini-0321:phyxminiproblems-problemphyxmini0321-matchesproblemandfigure}
  \lean{PhyXMiniProblems.ProblemPhyXMini0321.MatchesProblemAndFigure}
  \uses{def:physics:phyx-mini-0321:phyxminiproblems-problemphyxmini0321-durationinseconds, def:physics:phyx-mini-0321:phyxminiproblems-problemphyxmini0321-speedinmeterspersecond, def:physics:phyx-mini-0321:phyxminiproblems-problemphyxmini0321-spermwhaleclicksetup}
  The stated numerical data and the readouts from the primary figure. The 1.0 ms scale bar spans one third of the interval between corresponding click peaks. This predicate fixes no value for the requested sac length or for the round-trip distance
\end{definition}
\begin{proof}
  This is the declared model definition of Matches Problem And Figure.
\end{proof}
\begin{definition}[Has Physical Parameters]
  \label{def:physics:phyx-mini-0321:phyxminiproblems-problemphyxmini0321-hasphysicalparameters}
  \lean{PhyXMiniProblems.ProblemPhyXMini0321.HasPhysicalParameters}
  \uses{def:physics:phyx-mini-0321:phyxminiproblems-problemphyxmini0321-lengthinmeters, def:physics:phyx-mini-0321:phyxminiproblems-problemphyxmini0321-durationinseconds, def:physics:phyx-mini-0321:phyxminiproblems-problemphyxmini0321-speedinmeterspersecond, def:physics:phyx-mini-0321:phyxminiproblems-problemphyxmini0321-spermwhaleclicksetup}
  Strict positivity conditions for the physical branch of the model
\end{definition}
\begin{proof}
  This is the declared model definition of Has Physical Parameters.
\end{proof}
\begin{definition}[Satisfies Echo Round Trip Laws]
  \label{def:physics:phyx-mini-0321:phyxminiproblems-problemphyxmini0321-satisfiesechoroundtriplaws}
  \lean{PhyXMiniProblems.ProblemPhyXMini0321.SatisfiesEchoRoundTripLaws}
  \uses{def:physics:phyx-mini-0321:phyxminiproblems-problemphyxmini0321-lengthreadout, def:physics:phyx-mini-0321:phyxminiproblems-problemphyxmini0321-durationreadout, def:physics:phyx-mini-0321:phyxminiproblems-problemphyxmini0321-speedreadout, def:physics:phyx-mini-0321:phyxminiproblems-problemphyxmini0321-spermwhaleclicksetup}
  The governing geometry and acoustics laws. The echo route has two legs, each the length of the spermaceti sac, and propagation at constant speed obeys distance = speed * time. Both relations are required in every compatible choice of length and time units; neither contains the requested numerical answer
\end{definition}
\begin{proof}
  This is the declared model definition of Satisfies Echo Round Trip Laws.
\end{proof}
\begin{lemma}[spermaceti Sac Length formula]
  \label{lem:physics:phyx-mini-0321:phyxminiproblems-problemphyxmini0321-spermacetisaclength-formula}
  \lean{PhyXMiniProblems.ProblemPhyXMini0321.spermacetiSacLength_formula}
  \uses{def:physics:phyx-mini-0321:phyxminiproblems-problemphyxmini0321-lengthinmeters, def:physics:phyx-mini-0321:phyxminiproblems-problemphyxmini0321-durationinseconds, def:physics:phyx-mini-0321:phyxminiproblems-problemphyxmini0321-speedinmeterspersecond, def:physics:phyx-mini-0321:phyxminiproblems-problemphyxmini0321-spermwhaleclicksetup, def:physics:phyx-mini-0321:phyxminiproblems-problemphyxmini0321-hasphysicalparameters, def:physics:phyx-mini-0321:phyxminiproblems-problemphyxmini0321-satisfiesechoroundtriplaws}
  The general round-trip formula, before inserting this figure's numbers
\end{lemma}
\begin{proof}
  The conclusion follows from the governing relations and figure readouts named in the statement of spermaceti Sac Length formula.
\end{proof}
\begin{definition}[Answer Choice]
  \label{def:physics:phyx-mini-0321:phyxminiproblems-problemphyxmini0321-answerchoice}
  \lean{PhyXMiniProblems.ProblemPhyXMini0321.AnswerChoice}
  Labels of the four lengths displayed with the problem
\end{definition}
\begin{proof}
  This is the declared model definition of Answer Choice.
\end{proof}
\begin{definition}[meters]
  \label{def:physics:phyx-mini-0321:phyxminiproblems-problemphyxmini0321-answerchoice-meters}
  \lean{PhyXMiniProblems.ProblemPhyXMini0321.AnswerChoice.meters}
  \uses{def:physics:phyx-mini-0321:phyxminiproblems-problemphyxmini0321-answerchoice}
  Meter readout printed beside each displayed answer label
\end{definition}
\begin{proof}
  This is the declared model definition of meters.
\end{proof}
\begin{definition}[Matches Answer Choice]
  \label{def:physics:phyx-mini-0321:phyxminiproblems-problemphyxmini0321-matchesanswerchoice}
  \lean{PhyXMiniProblems.ProblemPhyXMini0321.MatchesAnswerChoice}
  \uses{def:physics:phyx-mini-0321:phyxminiproblems-problemphyxmini0321-acousticlength, def:physics:phyx-mini-0321:phyxminiproblems-problemphyxmini0321-lengthinmeters, def:physics:phyx-mini-0321:phyxminiproblems-problemphyxmini0321-answerchoice}
  A physical length agrees with a one-decimal-place displayed answer. The 0.05 m tolerance is half of the displayed 0.1 m increment
\end{definition}
\begin{proof}
  This is the declared model definition of Matches Answer Choice.
\end{proof}
\begin{definition}[recorded Answer Choice]
  \label{def:physics:phyx-mini-0321:phyxminiproblems-problemphyxmini0321-recordedanswerchoice}
  \lean{PhyXMiniProblems.ProblemPhyXMini0321.recordedAnswerChoice}
  \uses{def:physics:phyx-mini-0321:phyxminiproblems-problemphyxmini0321-answerchoice}
  The answer label recorded by the supplied dataset
\end{definition}
\begin{proof}
  This is the declared model definition of recorded Answer Choice.
\end{proof}
% --- Archon declaration topology end ---
% --- Archon physics formalization source end ---
